%
% Die Stammfunktion von 1/theta
%
\subsection{Die Stammfunktion von $1/\vartheta$
\label{buch:intalg:wurzel:subsection:stammfunktiontheta}}
Die Stammfunktion von $1/\vartheta$, das erste Integral in 
\eqref{buch:intalg:wurzel:eqn:grundintegrale},
hängt davon ab, wieviele reelle Nullstellen der Radikand hat und
ob $a$ positiv oder negativ ist.
Die Anzahl der Nullstellen ist durch das Vorzeichen der Diskriminante
$\Delta=b^2-4ac$ bestimmt.

\subsubsection{Der Fall $b^2-4ac=0$}
In diesem Fall hat der Radikand genau eine doppelte Nullstelle
bei $-b/2a$.
Für $a<0$ ist der Radikand nicht positiv, wir betrachten daher nur
den Fall $a>0$.
Die Substitution
\begin{align*}
x
&=
\frac{z}{\!\sqrt{a}}
z-\frac{b}{2a}
\qquad
\text{mit}
\qquad
dx
=
\frac{1}{\!\sqrt{a}}
dz
\\
z &= \!\sqrt{a}\biggl(x+\frac{b}{2a}\biggr)
\intertext{verwandelt den Integranden in}
ax^2+bx+c&=z^2
\intertext{und damit das Integral in}
\int\frac{1}{\vartheta}\,dx
&=
\int \frac{1}{\sqrt{z^2}} \frac{1}{\!\sqrt{a}}\,dz
=
\frac{1}{\!\sqrt{a}} \int\frac{1}{z}\,dz
=
\frac{1}{\!\sqrt{a}} \log z
\\
&=
\frac{1}{\!\sqrt{a}} \log \!\sqrt{a}\biggl(x+\frac{b}{2a}\biggr)
=
\frac{1}{\!\sqrt{a}} \log(2ax+b)
+
\frac{1}{\!\sqrt{a}} \log\frac{\sqrt{a}}{2a}
\intertext{Da die Stammfunktion nur bis auf eine Konstante bestimmt ist,
können wir den letzten Term in der Integrationskonstante absorbieren
und}
\int\frac{1}{\!\sqrt{ax^2+bx+c}}\,dx
&=
\frac{1}{\!\sqrt{a}} \log(2ax+b) + C
\end{align*}
schreiben.

\subsubsection{Der Fall $b^2-4ac>0$}
Für $b^2-4ac>0$ kann man durch die Substitution
\begin{align*}
x
&=
-\frac{b}{2a} + z\frac{\!\sqrt{b^2-4ac}}{2a}
&&\text{und}&
dx &= \frac{\!\sqrt{b^2-4ac}}{2a} \,dz
\\
z
&=
\frac{2a}{\!\sqrt{b^2-4ac}}\biggl(x+\frac{b}{2a}\biggr)
=
\frac{1}{\!\sqrt{b^2-4ac}} (2ax+b)
\end{align*}
den Radikanden in
\[
ax^2+bx+z
=
\frac{b^2-4ac}{4a}(1-z^2)
\]
umformen.
Für $a>0$ ist also
\begin{align*}
\int \frac{1}{\vartheta}\,dx
&=
\int \!\sqrt{\frac{4a}{b^2-4ac}} \frac{1}{\!\sqrt{1-z^2}}
\frac{\!\sqrt{b^2-4ac}}{2a} \,dz
\\
&=
\frac{1}{\!\sqrt{a}}
\int
\frac{1}{\!\sqrt{1-z^2}}\,dz
=
\frac{1}{\!\sqrt{a}}
\arcsin z
+C
\\
&=
\frac{1}{\!\sqrt{a}}
\arcsin
\frac{2ax+b}{\!\sqrt{b^2-4ac}}
+ C.
%
\intertext{Für $a<0$ ist der Integrand}
\vartheta
&=
\!\sqrt{\frac{b^2-4ac}{-4a}}
\!\sqrt{z^2-1}
\qquad\text{bzw.}\qquad
\!\sqrt{z^2-1}
=
\sqrt{\frac{-4a}{b^2-4ac}}\vartheta
\intertext{und damit das Integral in}
\int \frac{1}{\vartheta}\,dx
&=
\int \!\sqrt{\frac{-4a}{b^2-4ac}}\frac{1}{\!\sqrt{z^2-1}}
\frac{\!\sqrt{b^2-4ac}}{2a} \,dz
\\
&=
\frac{1}{\!\sqrt{-a}}
\int\frac{1}{\!\sqrt{z^2-1}}\,dz
=
\frac{1}{\!\sqrt{-a}}
\log(2\!\sqrt{z^2-1}+2z)
\\
&=
\frac{1}{\!\sqrt{-a}}
\log\biggl(
2\!\sqrt{\frac{-4a}{b^2-4ac}}\vartheta
+\frac{2}{\!\sqrt{b^2-4ac}}(2ax+b)
\biggr)
\\
&=
\frac{1}{\!\sqrt{-a}}
\log\bigl(
2\!\sqrt{-a}\vartheta
+
2ax+b
\bigr)
+
\frac{1}{\!\sqrt{-a}}
\log\biggl(\frac{2}{\!\sqrt{b^2-4ac}}\biggr).
\intertext{Der zweite Term ist eine Konstante, auf die es für die
Stammfunktion nicht ankommt.
Wir absorbieren ihn daher in der Integrationskonstanten und können
schreiben}
\int\frac{1}{\vartheta}\,dx
&=
\frac{1}{\!\sqrt{-a}}\log\bigl(2\sqrt{-a}\vartheta + 2ax+b\bigr)
+C.
\end{align*}

\subsubsection{Der Fall $b^2-4ac<0$}
In diesem Fall hat der Radikand keine Nullstelle.
Für $a<0$ ist der Radikand immer negativ, wir betrachten daher wieder
nur den Fall $a>0$.
Die Substitution
\begin{align*}
x&=-\frac{b}{2a}+z\frac{\!\sqrt{4ac-b^2}}{2a}
\qquad\text{und}\qquad
dx=\frac{4ac-b^2}{2a}\,dz
\\
z&=\frac{2a}{\!\sqrt{4ac-b^2}}\biggl(x+\frac{b}{2a}\biggr)
\intertext{verwandelt den Radikanden in}
ax^2+bx+c
&=
\frac{4ac-b^2}{4a}(z^2+1)
\intertext{und damit das Integral}
\int\frac{1}{\vartheta}\,dx
&=
\int
\!\sqrt{\frac{4a}{4ac-b^2}}
\frac{1}{ \!\sqrt{z^2+1}}
\frac{\!\sqrt{4ac-b^2}}{2a}
\,dz
=
\frac{1}{2\!\sqrt{a}}
\int \frac{1}{\!\sqrt{z^2+1}}\,dz
\\
&=
\frac{1}{\!\sqrt{a}}
\operatorname{arsinh} z+C
\\
&=
\frac{1}{\!\sqrt{a}}
\operatorname{arsinh} 
\frac{2a}{\!\sqrt{4ac-b^2}}\biggl(x+\frac{b}{2a}\biggr)+C
=
\frac{1}{\!\sqrt{a}}
\operatorname{arsinh} 
\frac{2ax+b}{\!\sqrt{4ac-b^2}}+C.
\end{align*}

%
% Reelle Ausdrücke für die Stammfunktion von 1/vartheta
%
\subsubsection{Reelle Ausdrücke für die Stammfunktion von $1/\vartheta$}
Wir fassen die fallweisen Rechnungen der vorangegangenen Abschnitt
ein einem Satz zusammen.

\begin{satz}
\label{buch:intalg:wurzel:satz:1theta}
Für $\vartheta = ax^2+bx+c$ mit reellen Koeffizienten $a,b,c\in\mathbb{R}$
gilt
\begin{equation}
\int\frac{1}{\vartheta}\,dx
=
\begin{cases}
\displaystyle
\frac{1}{\!\sqrt{a}}
\log(2ax+b)+C
&\text{falls $b^2-4ac=0$ und $a>0$,}
\\[10pt]
\displaystyle
\frac{1}{\!\sqrt{a}}
\arcsin\frac{2ax+b}{\!\sqrt{b^2-4ac}}+C
&\text{falls $b^2-4ac>0$ und $a>0$,}
\\[10pt]
\displaystyle
\frac{1}{\!\sqrt{-a}}
\log\bigl(2\!\sqrt{-a}\vartheta+2ax+b\bigr) + C
&\text{falls $b^2-4ac>0$ und $a<0$,}
\\[10pt]
\displaystyle
\frac{1}{\!\sqrt{a}}
\operatorname{arsinh}\frac{2ax+b}{\!\sqrt{4ac-b^2}}+C
&\text{falls $b^2-4ac<0$ und $a>0$.}
\end{cases}
\end{equation}
\end{satz}

%
% Ein gemeinsamer Ausdruck für alle Fälle
%
\subsubsection{Ein gemeinsamer Ausdruck für alle Fälle}
Die Fallunterscheidung von Satz~\ref{buch:intalg:wurzel:satz:1theta}
ist vor allem deshalb notwendig, weil wir imaginäre Werte von
$\!\sqrt{b^2-4ac}$ und $\!\sqrt{-a}$ vermeiden möchten.
Die Funktion $\arcsin$ und $\operatorname{arsinh}$ können aber
wie in Abschnitt~\ref{buch:intalg:elementar:subsection:trigo}
gezeigt auch ausschliesslich durch Logarithmen ausgedrückt werden.
Es sollte daher möglich sein, einen einzigen komplexen
Ausdruck für eine Stammfunktion von $1/\vartheta$ zu finden.

Durch Nachrechnen der Ableitung von
\[
\Theta
=
\frac{1}{\!\sqrt{a}}
\log\bigl(
2ax + b - 2\!\sqrt{a} \vartheta
\bigr)
\]
bekommt man
\begin{align*}
\Theta'
&=
-\frac{1}{\!\sqrt{a}}
\cdot
\frac{1}{2ax+b-2\!\sqrt{a}\vartheta}
\cdot
(2a - 2\!\sqrt{a}\vartheta')
\\
&=
-\frac{1}{\!\sqrt{a}}
\cdot
\frac{1}{2ax+b-2\!\sqrt{a}\vartheta}
\cdot
\biggl(
2a -2\!\sqrt{a}\frac{2ax+b}{2\vartheta}
\biggr)
\\
&=
-\frac{1}{\!\sqrt{a}}
\cdot
\frac{1}{2ax+b-2\!\sqrt{a}\vartheta}
\cdot
\frac{
4a\vartheta -2\!\sqrt{a}(2ax+b)
}{
2\vartheta
}
\\
&=
-\frac{1}{\!\sqrt{a}}
\cdot
\frac{1}{2ax+b-2\!\sqrt{a}\vartheta}
\cdot
\frac{
2\!\sqrt{a}\cdot
\bigl(
\!\sqrt{a}\vartheta
-
2ax-b
\bigr)
}{
2\vartheta
}
=
\frac{1}{\vartheta},
\end{align*}
Die Funktion $\Theta$ ist also eine Stammfunktion von $1/\vartheta$.
Sie liegt in einem Erweiterungskörper, dem die Konstante $\!\sqrt{a}$
und der Logarithmus $\log(2ax+b-2\!\sqrt{a}\vartheta)$ hinzugefügt
werden ist.
Damit der Logarithmus hinzugefügt werden kann, muss auch noch $b$
hinzugefügt werden.
Somit ist
\[
\Theta
\in
\mathbb{Q}\bigl(x,\theta,\!\sqrt{a},b,\log(\!\sqrt{a}\vartheta-2ax-b)\bigr)
\]
eine elementare Stammfunktion von $1/\vartheta$.

